\documentclass[11pt]{article}
\usepackage[margin=25mm]{geometry}
\usepackage{fontspec}
\setmainfont{Latin Modern Roman}
\usepackage{amsmath,amssymb,mathtools}
\usepackage{unicode-math}
\setmathfont{Latin Modern Math}
\usepackage{xurl,hyperref}
\hypersetup{colorlinks=true,urlcolor=blue}
\setcounter{secnumdepth}{0}
\setlength{\parindent}{0pt}
\setlength{\parskip}{5pt}
\setlength{\emergencystretch}{3em}
\newcommand{\sref}[2]{\hyperlink{source:#1:#2}{[S#2]}}
\newcommand{\cataloguescope}[1]{\paragraph{Recorded target.}#1}
\newcommand{\R}{\mathbb R}
\newcommand{\C}{\mathbb C}
\newcommand{\N}{\mathbb N}
\newcommand{\Z}{\mathbb Z}
\newcommand{\Q}{\mathbb Q}
\newcommand{\F}{\mathbb F}
\newcommand{\eps}{\varepsilon}
\begin{document}
% UnsolvedMath ID 30001175; source catalogue code OWR-3392-001
\setcounter{section}{30001174}
\section{Homogeneous Factorizations Producing Type-II Arveson Systems}\label{30001175}
{\small UnsolvedMath ID 30001175 · Mathematical Physics}
\subsection{Definitions and mathematical statement}

\paragraph{Shared notation.}$\mathbb N=\{1,2,\ldots\}$; the standard number fields have their usual meanings. All additional conventions are stated locally.


\paragraph{Definitions and conventions.} A factor is a von Neumann algebra with center $\mathbb C I$; type $\mathrm{II}_1$ is finite with a normalized trace and type $\mathrm{III}$ has no nonzero finite projections. The source uses a separable predual and a faithful normal state $\phi$. A noncommutative factorization assigns $\phi$-conditioned subalgebras $A_B$ to measurable subsets $B\subseteq\mathbb R$, respecting intersections, generated unions, independence of disjoint sets, and exhaustion. Homogeneity supplies state-preserving translations sending $A_B$ to $A_{B+t}$. An Arveson system is a measurable family of Hilbert spaces $E_t$, $t>0$, with associative unitaries $E_s\otimes E_t\cong E_{s+t}$. Type $\mathrm{II}$ means that units exist but do not generate the whole system.

\paragraph{Source scope and qualifications.} The report's factorization axioms, state-preserving conditional expectations, and continuity/measurability conditions are part of the construction; a family of arbitrary independent subalgebras is insufficient. The existence question concerns the given probability algebra, not simply existence of an unrelated type II system.

\paragraph{Catalogue mathematical target.}
\paragraph{Catalogue wording retained for traceability.}
If $A$ is a type $\mathrm{II}_1$ or type $\mathrm{III}$ factor with state $\phi$, does there exist a homogeneous noncommutative factorization of $(A,\phi)$ over $\mathbb R$ that leads to an Arveson system of type $\mathrm{II}$? Here homogeneity means a $\phi$-preserving translation action on the factorization over measurable subsets of $\mathbb R$.

\paragraph{Saved literature context.} The supplied catalogue's assessment dated 2026-08-21 records: The report poses existence of a homogeneous noncommutative factorization producing a type II Arveson system; no later construction resolving this general existence question was located. The assessment is retained as dated provenance; it is not a new comprehensive literature or open-status certification.
\subsection{Short English statement}
Construct a translation-invariant noncommutative factorization of the prescribed factor and state whose associated product system has type II.
\subsection{Research attempt}
\paragraph{Model, date and investigated scope.}
OpenAI GPT-6 Astra (configured model \texttt{gpt-6-astra}), ultra reasoning; generated 11 October 2026. Proves directly that the naive Gaussian Fock factorization is unit-generated and therefore cannot furnish the requested type-II construction.

\paragraph{Sources and literature check.}
The complete supplied definition, including source qualifications and the dated literature assessment, was reviewed. The calculation below is an elementary restricted attempt motivated by [S2]; it does not claim novelty or a fresh literature-status determination.

\paragraph{Context and investigated scope.} The complete supplied definition, including source qualifications and the dated literature assessment, was reviewed. The calculation below is an elementary restricted attempt motivated by [S2]; it does not claim novelty or a fresh literature-status determination.

\paragraph{Why the ordinary Gaussian construction gives the wrong type.} Let $K_t=L^2((0,t);\mathbb C)$ and $E_t=\bigoplus_{n\geq0}K_t^{\otimes_s n}$ be symmetric Fock space. For $f\in K_t$ define the exponential vector $e(f)=\bigoplus_{n\geq0}f^{\otimes n}/\sqrt{n!}$. Its squared norm is $e^{\|f\|^2}$, so it belongs to $E_t$. The orthogonal time decomposition identifies $e(f\oplus g)$ with $e(f)\otimes e(g)$.

\paragraph{Units already generate everything.} For each $z\in\mathbb C$, $u_t^z=e(z1_{(0,t)})$ is a measurable multiplicative section. Products of these units over finite time partitions give $e(f)$ for step functions $f$. Those vectors are total: if $v=(v_n)$ is orthogonal to every $e(f)$, then for any fixed finite collection of step functions the entire function $\langle v,e(\sum_jz_jf_j)\rangle$ is zero. Its Taylor coefficients imply orthogonality of each $v_n$ to all symmetric products of step functions. Step functions are dense in $K_t$, hence these products are dense in each symmetric tensor power, and every $v_n=0$.

\paragraph{Where the construction fails.} Therefore the ordinary Gaussian/Fock factorization is type I, not the requested type II. Adding more ordinary independent Gaussian channels repeats the same density argument. The source additionally prescribes a particular factor $(A,\phi)$; this Hilbert-space construction does not realize its state or conditional expectations. The first gap is a homogeneous factorization on that prescribed probability algebra with a nonzero component orthogonal to the subsystem generated by all units, while retaining at least one unit and all continuity axioms.

\paragraph{Verification and adversarial audit.}
Checked exponential-vector norms, step-function products and totality via Taylor coefficients; separated Hilbert-space type from prescribed algebra realization. Model self-audit only; no independent human proof checking. Prompts A and B were applied to this canonical record and its qualified source scope. Checked exponential-vector norms, step-function products and totality via Taylor coefficients; separated Hilbert-space type from prescribed algebra realization. Retained unresolved status and zero original-target progress; A homogeneous factorization of the prescribed factor/state with a proper unit-generated subsystem is missing.

Coordinator checked norms and factorization of exponential vectors and totality of the subsystem generated by step-function units. This additional check was model-assisted, not independent human verification.

\paragraph{Tools, process and contributions.}
Python 3 was used for file preparation and structural checks; no numerical experiment or theorem prover was used. Contribution prompt A was followed by prompt B. GPT-6 Astra with ultra reasoning drafted the research and performed its stated self-audit; any additional coordinator contribution is identified in the verification and audit. A GPT-6-based Codex coordinator assembled the manuscript and checked repository integration; its exact variant and reasoning configuration were not exposed. No human mathematical contribution, independent human proof check or formal proof certification is supplied. Complete derivations and the consulted source references are included here; private reasoning transcripts are not published.

\paragraph{Final status and first remaining gap.}
\textbf{UNRESOLVED}. A homogeneous factorization of the prescribed factor/state with a proper unit-generated subsystem is missing.

\section{Completion Estimate}
\noindent\textbf{COMPLETION: 0\%}
The calculation is an auxiliary or restricted attempt and establishes no newly verified portion of the full source target. A homogeneous factorization of the prescribed factor/state with a proper unit-generated subsystem is missing.

\subsection{Sources}
\begin{itemize}
\item[\hypertarget{source:30001175:1}{S1}] ULAM.AI and the UnsolvedMath Contributors, UnsolvedMath, record 30001175 (OWR-3392-001). Catalogue curation under CC BY 4.0; definitions and English exposition added locally.
\url{https://huggingface.co/datasets/ulamai/UnsolvedMath}.
\item[\hypertarget{source:30001175:2}{S2}] Mini-Workshop: Product Systems and Independence in Quantum Dynamics (2009), Oberwolfach Reports, DOI 10.4171/owr/2009/09, pp. 493--548; official report source and mathematical context.
\url{https://doi.org/10.4171/owr/2009/09}.
\item[\hypertarget{source:30001175:3}{S3}] Official EMS Press PDF of the same Oberwolfach report; PDF pages 29--30, printed pages 521--522.
\url{https://ems.press/content/serial-article-files/46209}.
\item[\hypertarget{source:30001175:4}{S4}] Tsirelson nonclassical stochastic flows; reference in the saved dated catalogue literature review.
\url{https://arxiv.org/abs/math/0402431}.
\end{itemize}
\label{end:30001175}
\end{document}
